\subsection{The instruments of explicit evaluation: a survey}
\label{p11:sec:ev-instruments}

The positivist evaluation of flourishing and of relationships has produced, over the past half-century, a substantial body of validated instruments, and the cultivation of a shared life that would assess itself explicitly draws upon them. This subsection surveys the principal instruments, classical and contemporary, and then locates them within the framework of the paper, in order to make precise both what they supply and what, by their construction, they cannot reach.

\subsubsection{The principal instruments and what they measure}
\label{p11:sec:ev-instr-survey}

The instruments divide by the construct they target. A first family measures \emph{subjective well-being} in the hedonic tradition, which construes well-being as life satisfaction together with the balance of positive over negative affect. The Satisfaction with Life Scale assesses the cognitive, evaluative component with five items \parencite{diener1985}; the Positive and Negative Affect Schedule assesses the affective component with twenty \parencite{watson1988}; and the tradition is gathered in \textcite{kahneman1999}. A second family measures \emph{eudaimonic well-being}, which construes well-being as the realisation of human potentials rather than the balance of feeling. Ryff's Scales of Psychological Well-Being operationalise this in six dimensions, autonomy, environmental mastery, personal growth, positive relations with others, purpose in life, and self-acceptance, in forms of eighteen, forty-two, and eighty-four items \parencite{ryff1989,ryffkeyes1995}; self-determination theory supplies measures grounded in the satisfaction of the basic needs for autonomy, competence, and relatedness \parencite{ryan2000,deci1985sdt}; and the PERMA framework adds engagement, relationships, meaning, and accomplishment to positive emotion \parencite{seligman2011}. A third family measures \emph{quality of life} through objective and mixed indicators, of which the World Health Organization's quality-of-life assessment, covering physical, psychological, social, and environmental domains, is representative \parencite{whoqol1998}. A fourth family, central to the present paper, measures \emph{relationship quality}: the Dyadic Adjustment Scale assesses the quality of a couple's relationship with thirty-two items across consensus, satisfaction, cohesion, and affectional expression \parencite{spanier1976}, while the Relationship Assessment Scale offers a brief, seven-item, unidimensional measure of global relationship satisfaction \parencite{hendrick1988}.

Contemporary work has extended this measurement from the individual and the couple to the population and the globe, and from a single index to multidimensional batteries. National and international instruments now assess well-being at the level of whole societies. The World Happiness Report ranks more than one hundred and forty countries on a single life-evaluation question, the Cantril ladder, on which respondents place their lives between the worst (0) and the best (10) possible, averaged over a three-year window from the Gallup World Poll \parencite{whr2026}. The OECD's Better Life Index assesses well-being across eleven dimensions, among them housing, income, jobs, community, education, environment, health, life satisfaction, safety, and work-life balance, drawn from twenty-four indicators and aggregated under user-assigned weights \parencite{oecdbli}. And the Global Flourishing Study, a longitudinal panel of more than two hundred thousand participants across twenty-two countries, measures flourishing across six domains, happiness and life satisfaction, physical and mental health, meaning and purpose, character and virtue, close social relationships, and financial and material stability \parencite{vanderweele2025gfs,vanderweele2017}. These instruments are summarised in Table~\ref{p11:tab:instruments}.

\begin{table}[t]
\centering
\small
\renewcommand{\arraystretch}{1.25}
\caption{Principal instruments for the quantitative measurement of well-being, quality of life, and relationship quality.}
\label{p11:tab:instruments}
\begin{tabularx}{\linewidth}{@{}>{\raggedright\arraybackslash}p{3.0cm} >{\raggedright\arraybackslash}p{2.4cm} >{\raggedright\arraybackslash}X >{\raggedright\arraybackslash}p{2.2cm}@{}}
\toprule
\rowcolor{pageblush}
\textbf{\textcolor{rosedeep}{Instrument}} & \textbf{\textcolor{rosedeep}{Tradition}} & \textbf{\textcolor{rosedeep}{What it measures}} & \textbf{\textcolor{rosedeep}{Source}} \\
\midrule
Satisfaction with Life Scale & Hedonic SWB & Cognitive life satisfaction; 5 items, 7-point & \textcite{diener1985} \\
PANAS & Hedonic SWB & Positive and negative affect; 20 items & \textcite{watson1988} \\
Ryff Scales (PWB) & Eudaimonic & Six dimensions of positive functioning; 18--84 items & \textcite{ryff1989,ryffkeyes1995} \\
Self-determination measures & Eudaimonic & Autonomy, competence, relatedness need-satisfaction & \textcite{ryan2000} \\
PERMA & Eudaimonic & Positive emotion, engagement, relationships, meaning, accomplishment & \textcite{seligman2011} \\
WHOQOL & Quality of life & Physical, psychological, social, environmental domains & \textcite{whoqol1998} \\
Dyadic Adjustment Scale & Relationship & Consensus, satisfaction, cohesion, affection; 32 items & \textcite{spanier1976} \\
Relationship Assessment Scale & Relationship & Global relationship satisfaction; 7 items, unidimensional & \textcite{hendrick1988} \\
World Happiness Report & Macro/national & Single life-evaluation (Cantril ladder); national averages & \textcite{whr2026} \\
OECD Better Life Index & Macro/national & Eleven dimensions, 24 indicators, user-weighted & \textcite{oecdbli} \\
Global Flourishing Study & Eudaimonic/macro & Six domains incl.\ character, meaning, relationships; longitudinal & \textcite{vanderweele2025gfs} \\
\bottomrule
\end{tabularx}
\end{table}

\subsubsection{The instruments within the framework of this paper}
\label{p11:sec:ev-instr-frame}

The survey makes possible a precise statement of what the existing instruments supply to the assessment of a shared flourishing and what they do not. Located within the framework of this paper, each instrument proxies some component of the normative specification of \xref{p11:sec:cog-spec}, and the location reveals a systematic gap. The hedonic instruments proxy the felt quality of a life, an aspect of the eudaimonic condition, and they proxy it at the level of the aggregate or the mean, which is exactly the level at which the maldistribution of return is invisible (\xref{p11:sec:ev-skew}). The eudaimonic instruments proxy the unfolding of each party's potentials, and so come nearest the eudaimonic condition of the unfolding of each, but they measure the individual's functioning in isolation and do not assess whether one party's unfolding is purchased by another's foreclosure. The relationship instruments proxy the parties' satisfaction with the relation, again at the level of each party's report, and a relation may score well on them while distributing its returns unjustly, since a favoured party's high satisfaction and an habituated party's adjusted satisfaction together yield an unremarkable aggregate. The macro instruments proxy well-being at the level of populations and are not addressed to the internal justice of a particular shared life at all.

The systematic gap is the condition of justice. No instrument in the survey measures whether the value a party generates returns to that party, nor whether one party is consigned to the low-return tail of \xref{p11:sec:ev-skew}; the instruments measure satisfaction, functioning, and conditions, each at the level of the party or the aggregate, and the distribution of return among the parties, relative to their coupling, is not among the constructs they target. This is not a defect of the instruments within their own traditions but a boundary of those traditions, and it is the boundary at which the paper's quantitative assessment of justice (\xref{p11:sec:ev-justice}) is directed. Table~\ref{p11:tab:framework} locates the instruments within the framework, by the component of a just eudaimonia each proxies and by the moment of the cycle at which each is principally drawn upon, and it records, in its final column, that the condition of justice is the one no existing family of instruments supplies.

\begin{table}[t]
\centering
\small
\renewcommand{\arraystretch}{1.25}
\caption{The instruments located within the framework of the paper: the component of a just eudaimonia each proxies, the moment of the cycle at which each is drawn upon, and whether it registers the condition of justice.}
\label{p11:tab:framework}
\begin{tabularx}{\linewidth}{@{}>{\raggedright\arraybackslash}p{2.7cm} >{\raggedright\arraybackslash}X >{\raggedright\arraybackslash}p{2.3cm} >{\raggedright\arraybackslash}p{2.3cm}@{}}
\toprule
\rowcolor{pageblush}
\textbf{\textcolor{rosedeep}{Family}} & \textbf{\textcolor{rosedeep}{Component of a just eudaimonia proxied}} & \textbf{\textcolor{rosedeep}{Cycle moment}} & \textbf{\textcolor{rosedeep}{Registers justice?}} \\
\midrule
Hedonic SWB & Felt quality of life, at the aggregate or mean & Evaluation & No; mean conceals skew \\
Eudaimonic & Unfolding of each party's potentials, in isolation & Cognition, evaluation & Partially; misses foreclosure by the other \\
Relationship quality & Each party's satisfaction with the relation & Evaluation, reflection & No; aggregate conceals maldistribution \\
Quality of life & Objective conditions of a life & Practice, evaluation & No; not addressed to return \\
Macro/national & Population-level well-being & (Not intra-relational) & No; not addressed to a particular relation \\
\rowcolor{pageblush}
Skew of return \emph{(this paper)} & Justice: return to each relative to coupling & Evaluation, reflection & Yes; this is its object \\
\bottomrule
\end{tabularx}
\end{table}

The two tables together state the contribution of the paper's quantitative assessment against the background of the existing instruments. The existing instruments, classical and contemporary, measure satisfaction, functioning, conditions, and their population aggregates, and they do so with a validity and an accumulability the paper does not match and does not attempt to match. What they do not measure, by the construction of the constructs they target, is the distribution of return among the parties of a relation, and so the condition of justice that the prior paper held to be constitutive of eudaimonia. The skew of the return distribution is offered as the structured proxy for exactly this gap: not a replacement for the existing instruments, but an addition orthogonal to them, addressed to the one component of a just eudaimonia that the survey shows no existing family of instruments to register.

\subsection{The operationalisation of return, and a worked illustration}
\label{p11:sec:ev-operational}

The assessment of justice requires that the abstract notion of the return of value be operationalised into a quantity that could, in principle, be estimated; this subsection sets out the operationalisation and illustrates it on a deliberately simple example, under the operational concession throughout.

The operationalisation proceeds in four steps. First, the contributions of each party over a period are identified: the discrete acts of generation, the pieces of sustaining labour, the cared-for occasions, that each party contributes to the shared life. Second, to each contribution is assigned a value, in a common unit, proxying the value it generates for the shared life; the assignment is the most contestable step, and the concession applies most stringently here, since the value of an act of care is exactly the unsymbolised quantity the proxy cannot reach, and the assignment is an operational stipulation and not a measurement. Third, for each party the return is identified: of the value generated in the shared life, the portion that returns to that party, whether directly or through the relation's circulation. Fourth, the return to each party is normalised by the coupling-based baseline of \xref{p11:sec:ev-return}, the return the coupling structure would yield under unbiased propagation, to give the normalised return on which the assessment of skewness is conducted.

A worked illustration, with stipulated numbers, shows the quantities in operation; the numbers are illustrative and carry no empirical claim. Consider a deeply coupled relation of two parties, $A$ and $B$, over a period in which the shared life generates value, in the common unit, of 100. Suppose $B$ bears the larger share of the sustaining labour and generates 60 of this value, while $A$ generates 40. In a just cycle, given the depth of the coupling, the return to each would be approximately proportional to generation: $B$ would receive a return near 60 and $A$ a return near 40, the normalised returns of both being near 1, and the distribution of normalised return approximately symmetric. Suppose instead that the relation returns value disproportionately to $A$: of the 100 generated, $A$ receives a return of 65 and $B$ a return of 35. The normalised return to $A$ is then $65/40 \approx 1.6$ and to $B$ is $35/60 \approx 0.58$. The distribution of normalised return is now skewed: $A$ sits well above the baseline of 1, and $B$ well below it, in the low-return tail. The skewness of the normalised-return distribution, negative for $B$, is the statistical fingerprint of \xref{p11:sec:ev-skew}, and it is legible here despite the relation's possibly high mean satisfaction, since $A$, the over-returned party, may report great contentment, and $B$, habituated to bearing the sustaining labour for little return, may report a satisfaction only moderately depressed, so that the mean of the two is unremarkable while the skew is severe. The illustration shows the quantities the assessment computes; it does not show that they can be estimated easily in an actual relation, which, as \xref{p11:sec:ev-model} stated, remains a substantial empirical problem the paper does not solve.
